%!TEX root = tfg.tex
\chapter{Introducción}

\begin{quotation}[Escritor]{Mark Twain (1835--1910)}
El secreto para avanzar es comenzar.
\end{quotation}

\begin{quotation}[Novelist]{Ernest Hemingway (1899--1961)}
The good parts of a book may be only something a writer is lucky enough to overhear or it may be the wreck of his whole damn life -- and one is as good as the other.
\end{quotation}

\begin{abstract}
Este capítulo presenta el punto de partida del proyecto: cómo se ha organizado el desarrollo, qué herramientas se han utilizado y cómo se ha documentado el proceso a lo largo del trabajo.
\end{abstract}

\section{Organización del estudio previo}

Todo proyecto de ingeniería requiere, antes de su desarrollo, una fase de planificación que defina cómo se va a organizar el trabajo, qué herramientas se van a utilizar y cómo se va a documentar el proceso. En este caso, la naturaleza del proyecto, que combina desarrollo de software en ROS2, experimentación con simulación y redacción de la memoria, ha condicionado las decisiones metodológicas adoptadas.

El estudio previo recoge las decisiones tomadas en cuanto a la organización del desarrollo, el control de versiones y la documentación del proceso, con el objetivo de que el trabajo sea reproducible y trazable.
